% xr's old release, which \usepackage{xr}[=v5] loads: its \externaldocument
% takes no second optional argument, and \externalcitedocument is a copy
% of it. The page builds, and each reference prints as in the PDF: a label
% of an .aux beside the document, with a prefix, a label of an .aux that
% LaTeX wrote with its number or page, one of an .aux that tex4ht wrote,
% with its .xref, as a link to that document's page, one with no prefix,
% one through \externalcitedocument, and one of a file that does not exist
% yet as ??. An \@input line of the other .aux, which \include writes, is
% read as written, beside this document, whatever folder the last
% \externaldocument names. So the \@input{chap.aux} of test-xr/latex.aux
% prints ??, as in the PDF.
\documentclass{article}
\usepackage{xr}[=v5]
\externaldocument[S-]{test-xrold-ext}
\externaldocument[X-]{test-xr/latex}
\externaldocument[T-]{test-xrold/t4ht}
\externaldocument{test-xrold/ext}
\externaldocument[M-]{test-xrold/none}
\externalcitedocument[C-]{test-xr/plain}
\newtheorem{theorem}{Theorem}
\begin{document}
CaseA \ref{S-five}, included \ref{S-six}.

CaseB \ref{X-ss}, page \pageref{X-front}, included \ref{X-c}.

CaseC \ref{T-oss}.

CaseD \ref{seven}, included \ref{eight}.

CaseE \ref{M-a}.

CaseF cited \ref{C-s}, page \pageref{C-front}.

\textit{In italics, \ref{X-ss}.}
A note.\footnote{In a footnote, \ref{X-ss}.}
\textsf{In sans-serif, \ref{X-ss}.}
\texttt{In typewriter, \ref{X-ss}.}

\begin{theorem}
In a theorem, \ref{X-ss}.
\end{theorem}

\section{See \ref{X-ss}}

\section{Here}\label{here}
Local \ref{here}.
\end{document}
